\documentclass{amsart}
\usepackage{amssymb,amsthm,mathtools}
\usepackage{hyperref}
\usepackage{cleveref}
\DeclarePairedDelimiter{\gen}{\langle}{\rangle}
\DeclarePairedDelimiter{\abs}{\lvert}{\rvert}
\DeclarePairedDelimiter{\floor}{\lfloor}{\rfloor}
\DeclarePairedDelimiter{\norm}{\lVert}{\rVert}
\DeclareMathOperator{\Hom}{Hom}
\DeclareMathOperator{\Aut}{Aut}
\newtheorem{thm}{Theorem}
\newtheorem{lem}[thm]{Lemma}
\theoremstyle{definition}
\newtheorem{defn}[thm]{Definition}
\theoremstyle{remark}
\newtheorem{rmk}[thm]{Remark}
\newtheorem*{thm*}{Theorem}
\crefname{thm}{Theorem}{Theorems}
\crefname{lem}{Lemma}{Lemmas}
\begin{document}
%% feature: paired-delims | \DeclarePairedDelimiter plain, starred (auto-size) and sized
\[ \abs{x}, \quad \abs*{\frac{a}{b}}, \quad \norm[\big]{v}, \quad \gen{g, h}, \quad \floor*{\frac{n}{2}} \]
%% feature: colon-eq | \coloneqq, \eqqcolon, \colon, \mapsto
\[ f \colon X \to Y, \quad x \mapsto x^2, \quad A \coloneqq B, \quad C \eqqcolon D \]
%% feature: ext-arrows | \xrightarrow with sub/superscript, \xhookrightarrow, \xleftrightarrow, \xRightarrow
\[ A \xrightarrow[\sim]{f} B \xhookrightarrow{i} C \xleftrightarrow{\cong} D \xRightarrow{\ \alpha\ } E \]
%% feature: align | align with numbered rows, \label/\eqref, and align*
\begin{align}
  (a+b)^2 &= a^2 + 2ab + b^2 \label{eq:sq} \\
  (a-b)^2 &= a^2 - 2ab + b^2
\end{align}
\begin{align*}
  \Hom(A, B) &\cong \Aut(B) & x &= y
\end{align*}
By \eqref{eq:sq} we are done.
%% feature: cases-matrices | cases, pmatrix, bmatrix, substack, binom, underbrace
\[ f(x) = \begin{cases} 1 & x > 0, \\ 0 & \text{otherwise}, \end{cases}
   \quad \begin{pmatrix} a & b \\ c & d \end{pmatrix}
   \begin{bmatrix} 1 \\ 0 \end{bmatrix},
   \quad \sum_{\substack{i < j \\ i,j \in S}} \binom{n}{k},
   \quad \underbrace{1 + \cdots + 1}_{n} \]
%% feature: big-operators | big operators with limits, \bigoplus, \bigsqcup, \prod, \int
\[ \bigoplus_{i \in I} M_i, \quad \bigsqcup_{n} X_n, \quad \prod_{p} (1 - p^{-s})^{-1}, \quad \int_{-\infty}^{\infty} e^{-x^2} \, dx \]
%% feature: theorems | plain/definition/remark styles, shared counter, unnumbered, proof with \qedhere
\begin{thm}\label{thm:main}
  Every finite group of prime order is cyclic.
\end{thm}
\begin{defn}
  A group is \emph{simple} if it has no nontrivial normal subgroups.
\end{defn}
\begin{rmk}
  This is a remark.
\end{rmk}
\begin{thm*}
  An unnumbered theorem.
\end{thm*}
\begin{proof}
  Lagrange's theorem gives
  \[ |G| = p. \qedhere \]
\end{proof}
%% feature: cleveref | \cref/\Cref with custom names to theorem and lemma
\begin{lem}\label{lem:a}
  A lemma.
\end{lem}
\begin{thm}\label{thm:b}
  A theorem.
\end{thm}
See \Cref{lem:a} and \cref{thm:b}, or both: \cref{lem:a,thm:b}.
\end{document}
